\subsection{Tensor product by a flat complex bounded on the right}

Lemma 1. --- Let C be a complex of right A-modules and E a complex of left A-modules. Suppose that H(C) = 0 and that E is flat and bounded on the right. Then H(C ⊗\_A E) = 0.

For \(k \in \mathbf{Z}\) , let \(\mathbf{T}^{(k)}\) be the subcomplex of \(\mathbf{C} \otimes_{\mathbf{A}} \mathbf{E}\) such that

\[
\mathrm{T}_{n}^{(k)} = \sum_{\substack{p + q = n\\ q\leqslant k}}\mathrm{C}_{p}\otimes_{\mathrm{A}}\mathrm{E}_{q};
\]

then \(\mathbf{T}^{(k - 1)}\subset \mathbf{T}^{(k)}\) and we have an exact sequence of complexes

\[
0 \longrightarrow \mathrm{T} ^ {(k - 1)} \stackrel {i _ {k}} {\longrightarrow} \mathrm{T} ^ {(k)} \stackrel {\pi} {\longrightarrow} \mathrm{C} \otimes_ {\mathbf {A}} \mathrm{E} _ {k} (- k) \longrightarrow 0
\]

where \(i_{k}\) is the canonical injection and where \(\pi\) projects the preceding direct sum onto its factor \(\mathbf{C}_{n-k} \otimes_{\mathbf{A}} \mathbf{E}_{k} = (\mathbf{C} \otimes_{\mathbf{A}} \mathbf{E}_{k}(-k))_{n}\) . By cor. 2 above, we have \(\mathbf{H}(\mathbf{C} \otimes_{\mathbf{A}} \mathbf{E}_{k}(-k)) = 0\) , hence \(i_{k}\) is a homologism. We have \(\mathbf{T}^{(k)} = 0\) for \(k\) small enough, since E is bounded on the right, hence \(\mathrm{H}(\mathrm{T}^{(k)}) = 0\) for all k by induction on k. Finally, the canonical morphism \(\lim_{\mathbf{A}} \mathrm{T}^{(k)} \to \mathbf{C} \otimes_{\mathbf{A}} \mathbf{E}\) is obviously an isomorphism, hence \(\mathrm{H}(\mathbf{C} \otimes_{\mathbf{A}} \mathbf{E}) = 0\) (X, p.~28, prop. 1).

Lemma 2. --- If \(u: \mathbf{C} \to \mathbf{C}'\) is a morphism of complexes of right A-modules and E is a complex of left A-modules, then the complexes Con \((u) \otimes_{\mathbf{A}} \mathbf{E}\) and Con \((u \otimes 1_{\mathbf{E}})\) are isomorphic.

By definition, Con \((u) \otimes_{\mathbf{A}} \mathbf{E}\) is the graded module \((\mathbf{C}'(-1) \oplus \mathbf{C}) \otimes_{\mathbf{A}} \mathbf{E}\) endowed with the differential D such that, for \(x \in \mathbf{C}_p\) , \(y' \in \mathbf{C}'(-1)_p = \mathbf{C}_{p-1}'\) , \(z \in \mathbf{E}_q\) , one has

\[
\mathrm{D} ((y ^ {\prime}, x) \otimes z) = (- d y ^ {\prime}, d x - u (y ^ {\prime})) \otimes z + (- 1) ^ {p} (y ^ {\prime}, x) \otimes d z,\tag{4}
\]

while \(\operatorname{Con}(u \otimes 1_{\mathbf{E}})\) is the graded module \((\mathbf{C}' \otimes_{\mathbf{A}} \mathbf{E})(-1) \oplus (\mathbf{C} \otimes_{\mathbf{A}} \mathbf{E})\) equipped with the differential \(\mathbf{D}_1\) such that, for \(x \in \mathbf{C}_p\) , \(y' \in \mathbf{C}_{p-1}'\) , \(z \in \mathbf{E}_q\) , one has

\[
\mathrm{D} _ {1} \left(y ^ {\prime} \otimes z, x \otimes z\right) = (- d y ^ {\prime} \otimes z - (- 1) ^ {p - 1} y ^ {\prime} \otimes d z, d x \otimes z + (- 1) ^ {p} x \otimes d z - u \left(y ^ {\prime}\right) \otimes z)
\]

\[
= \left(- d y ^ {\prime} \otimes z, (d x - u (y ^ {\prime})) \otimes z\right) + (- 1) ^ {p} \left(y ^ {\prime} \otimes d z, x \otimes d z\right)
\]

whence the assertion.

PROPOSITION 4. — Let \(u: \mathbf{C} \to \mathbf{C}'\) be a homologism of complexes of right A-modules, and E a complex of left A-modules, flat and bounded on the right. Then

\[
u \otimes 1 _ {\mathbf {E}}: \mathbf {C} \otimes_ {\mathbf {A}} \mathbf {E} \rightarrow \mathbf {C} ^ {\prime} \otimes_ {\mathbf {A}} \mathbf {E}
\]

is a homologism of complexes of k-modules.

Indeed, according to X, p.~38, cor., \(u\) (resp. \(u \otimes 1_{\mathrm{E}}\) ) is a homologism if and only if \(\mathrm{H}(\mathrm{Con}(u)) = 0\) (resp. \(\mathrm{H}(\mathrm{Con}(u \otimes 1_{\mathrm{E}})) = 0\) ). We then conclude by Lemmas 1 and 2.

Remark. --- Using the commutation isomorphisms, one deduces from the preceding statements the analogous statements obtained by interchanging the roles of the two arguments of the tensor products.

